\documentclass[10pt,a4paper]{amsart}
\usepackage{amsmath,amssymb,amsthm,graphicx}
% Retained minimal preamble for the Stacks Project chapter more-algebra.tex (More on Algebra).
% The chapter file pinned-original-source.tex carries only \input{preamble} on its first line, so no
% preamble text lives inside the chapter itself.  The Stacks Project uses one single shared preamble for
% every chapter; the snapshot stages it identically as sources/stacks-chow/preamble.tex and
% sources/stacks-etale-cohomology/preamble.tex.  The block below is copied mechanically, line for line,
% from sources/stacks-chow/preamble.tex lines 14-18 (the Xy-pic setup) and 154-217 (the native theorem
% declarations together with the mathematical macros).  Nothing was added, removed or redefined; the
% wrapper class and the amsmath/amssymb/amsthm/graphicx base packages are supplied by the assembler.
% Chapter: more-algebra.tex -- More on Algebra.
% For commutative diagrams we use Xy-pic
\usepackage[all]{xy}

% We use 2cell for 2-commutative diagrams.
\xyoption{2cell}

% Theorem environments.
%
\theoremstyle{plain}
\newtheorem{theorem}[subsection]{Theorem}
\newtheorem{proposition}[subsection]{Proposition}
\newtheorem{lemma}[subsection]{Lemma}

\theoremstyle{definition}
\newtheorem{definition}[subsection]{Definition}
\newtheorem{example}[subsection]{Example}
\newtheorem{exercise}[subsection]{Exercise}
\newtheorem{situation}[subsection]{Situation}

\theoremstyle{remark}
\newtheorem{remark}[subsection]{Remark}
\newtheorem{remarks}[subsection]{Remarks}

\numberwithin{equation}{subsection}

% Macros
%
\def\lim{\mathop{\mathrm{lim}}\nolimits}
\def\colim{\mathop{\mathrm{colim}}\nolimits}
\def\Spec{\mathop{\mathrm{Spec}}}
\def\Hom{\mathop{\mathrm{Hom}}\nolimits}
\def\Ext{\mathop{\mathrm{Ext}}\nolimits}
\def\SheafHom{\mathop{\mathcal{H}\!\mathit{om}}\nolimits}
\def\SheafExt{\mathop{\mathcal{E}\!\mathit{xt}}\nolimits}
\def\Sch{\mathit{Sch}}
\def\Mor{\mathop{\mathrm{Mor}}\nolimits}
\def\Ob{\mathop{\mathrm{Ob}}\nolimits}
\def\Sh{\mathop{\mathit{Sh}}\nolimits}
\def\NL{\mathop{N\!L}\nolimits}
\def\CH{\mathop{\mathrm{CH}}\nolimits}
\def\proetale{{pro\text{-}\acute{e}tale}}
\def\etale{{\acute{e}tale}}
\def\QCoh{\mathit{QCoh}}
\def\Ker{\mathop{\mathrm{Ker}}}
\def\Im{\mathop{\mathrm{Im}}}
\def\Coker{\mathop{\mathrm{Coker}}}
\def\Coim{\mathop{\mathrm{Coim}}}

% Boxtimes
%
\DeclareMathSymbol{\boxtimes}{\mathbin}{AMSa}{"02}

%
% Macros for moduli stacks/spaces
%
\def\QCohstack{\mathcal{QC}\!\mathit{oh}}
\def\Cohstack{\mathcal{C}\!\mathit{oh}}
\def\Spacesstack{\mathcal{S}\!\mathit{paces}}
\def\Quotfunctor{\mathrm{Quot}}
\def\Hilbfunctor{\mathrm{Hilb}}
\def\Curvesstack{\mathcal{C}\!\mathit{urves}}
\def\Polarizedstack{\mathcal{P}\!\mathit{olarized}}
\def\Complexesstack{\mathcal{C}\!\mathit{omplexes}}
% \Pic is the operator that assigns to X its picard group, usage \Pic(X)
% \Picardstack_{X/B} denotes the Picard stack of X over B
% \Picardfunctor_{X/B} denotes the Picard functor of X over B
\def\Pic{\mathop{\mathrm{Pic}}\nolimits}
\def\Picardstack{\mathcal{P}\!\mathit{ic}}
\def\Picardfunctor{\mathrm{Pic}}
\def\Deformationcategory{\mathcal{D}\!\mathit{ef}}

% --- appended 2026-10-12 (Chief): the project own definitions of the reference/slogan/history
% display environments, copied mechanically from the shared Stacks preamble lines 8-12
% (\usepackage{verbatim} plus the three \newenvironment lines, which the chapter preamble keeps
% verbatim from the project preamble). Nothing added, removed or redefined: these environments
% carry no mathematical content and the project itself typesets them as comments. Without these
% three lines a counted statement whose own text uses them fails with Environment slogan undefined.
\usepackage{verbatim}
\newenvironment{reference}{\comment}{\endcomment}
\newenvironment{slogan}{\comment}{\endcomment}
\newenvironment{history}{\comment}{\endcomment}
\usepackage{cleveref}
\crefname{originallemma}{original lemma}{original lemmas}
\Crefname{originallemma}{Original lemma}{Original lemmas}
\makeatletter
\newlabel{algebra-lemma-flat-going-down}{{00HS}{source}{}{sourcepointer-1}{}}
\newlabel{algebra-lemma-flat-going-down@cref}{{[originallemma][0][]00HS}{[1][0][]source}{}{}{}}
\makeatother
\crefname{originallemma}{original lemma}{original lemmas}
\Crefname{originallemma}{Original lemma}{Original lemmas}
\makeatletter
\newlabel{algebra-lemma-geometrically-normal-tensor-normal}{{06DF}{source}{}{sourcepointer-2}{}}
\newlabel{algebra-lemma-geometrically-normal-tensor-normal@cref}{{[originallemma][0][]06DF}{[1][0][]source}{}{}{}}
\makeatother
\crefname{originallemma}{original lemma}{original lemmas}
\Crefname{originallemma}{Original lemma}{Original lemmas}
\makeatletter
\newlabel{algebra-lemma-integral-going-up}{{00GU}{source}{}{sourcepointer-3}{}}
\newlabel{algebra-lemma-integral-going-up@cref}{{[originallemma][0][]00GU}{[1][0][]source}{}{}{}}
\makeatother
\crefname{originallemma}{original lemma}{original lemmas}
\Crefname{originallemma}{Original lemma}{Original lemmas}
\makeatletter
\newlabel{algebra-lemma-quotient-strict-henselization}{{05WS}{source}{}{sourcepointer-4}{}}
\newlabel{algebra-lemma-quotient-strict-henselization@cref}{{[originallemma][0][]05WS}{[1][0][]source}{}{}{}}
\makeatother
\crefname{originallemma}{original lemma}{original lemmas}
\Crefname{originallemma}{Original lemma}{Original lemmas}
\makeatletter
\newlabel{algebra-lemma-unique-prime-over-localize-below}{{00EA}{source}{}{sourcepointer-5}{}}
\newlabel{algebra-lemma-unique-prime-over-localize-below@cref}{{[originallemma][0][]00EA}{[1][0][]source}{}{}{}}
\makeatother
\crefname{originallemma}{original lemma}{original lemmas}
\Crefname{originallemma}{Original lemma}{Original lemmas}
\makeatletter
\newlabel{lemma-geometrically-unibranch}{{31453}{source}{}{sourcepointer-6}{}}
\newlabel{lemma-geometrically-unibranch@cref}{{[originallemma][0][]31453}{[1][0][]source}{}{}{}}
\makeatother
\title{The minimal primes of the strict henselization of a tensor product of strictly henselian local algebras correspond to pairs of minimal primes}
\author{The Stacks Project Authors}
\date{stacks-project:more-algebra.tex:local-snapshot-20261012; GNU FDL 1.2 or later}
\begin{document}
\maketitle
\setcounter{subsection}{1}
\noindent\textit{Editorial note.} All delivered mathematical text is the author's own, copied from the pinned Stacks Project chapter file more-algebra.tex; no statement or proof step was rewritten, reordered or removed.\par\medskip
\noindent\textit{Editorial note.} The chapter file pinned-original-source.tex carries only \textbackslash{}input\{preamble\} on its first line, so no preamble text lives inside the chapter. The retained file sources/stacks-more-algebra/retained-preamble.tex is a mechanical, line-for-line copy of the project's single shared preamble as staged in the snapshot; nothing was added, removed or redefined and no package was added.\par\medskip
\noindent\textit{Editorial note.} The source URL is the source-level Stacks Project URL, as required when more than one card is produced from one source; the per-claim tag is carried by claim\_key.\par\medskip
\noindent\textit{Editorial note.} The seven references inside the counted proof span are recorded as non-load-bearing navigation pointers; they name structural facts about strict henselization and minimal primes whose applications stay inside the counted proof span.\par\medskip
\noindent\textit{Allowed standard prerequisites (editorial).}\begin{itemize}
\item Strict henselization of a local ring as a filtered colimit of etale neighbourhoods and its residue field. used throughout the counted proof span. Is the construction whose minimal primes are being described.
\item Minimal primes of a tensor product of algebras over an algebraically closed field. used in the counted proof span. Supplies the pair-of-minimal-primes description that is transported to C.
\end{itemize}\medskip
\ifdefined\hypertarget\hypertarget{sourcepointer-1}{}\fi
\noindent\textit{Source reference.} Original lemma 00HS: focused-seeds/AGP/sources/stacks-snapshot-20261011/tags.txt, source lines 654--654. Cross-chapter navigation pointer into the Stacks Project; the standard fact it supplies is applied inside the counted proof span and the argument resting on it stays there.\par\medskip
\ifdefined\hypertarget\hypertarget{sourcepointer-2}{}\fi
\noindent\textit{Source reference.} Original lemma 06DF: focused-seeds/AGP/sources/stacks-snapshot-20261011/tags.txt, source lines 8179--8179. Cross-chapter navigation pointer into the Stacks Project; the standard fact it supplies is applied inside the counted proof span and the argument resting on it stays there.\par\medskip
\ifdefined\hypertarget\hypertarget{sourcepointer-3}{}\fi
\noindent\textit{Source reference.} Original lemma 00GU: focused-seeds/AGP/sources/stacks-snapshot-20261011/tags.txt, source lines 620--620. Cross-chapter navigation pointer into the Stacks Project; the standard fact it supplies is applied inside the counted proof span and the argument resting on it stays there.\par\medskip
\ifdefined\hypertarget\hypertarget{sourcepointer-4}{}\fi
\noindent\textit{Source reference.} Original lemma 05WS: focused-seeds/AGP/sources/stacks-snapshot-20261011/tags.txt, source lines 7596--7596. Cross-chapter navigation pointer into the Stacks Project; the standard fact it supplies is applied inside the counted proof span and the argument resting on it stays there.\par\medskip
\ifdefined\hypertarget\hypertarget{sourcepointer-5}{}\fi
\noindent\textit{Source reference.} Original lemma 00EA: focused-seeds/AGP/sources/stacks-snapshot-20261011/tags.txt, source lines 527--527. Cross-chapter navigation pointer into the Stacks Project; the standard fact it supplies is applied inside the counted proof span and the argument resting on it stays there.\par\medskip
\ifdefined\hypertarget\hypertarget{sourcepointer-6}{}\fi
\noindent\textit{Source reference.} Original lemma 31453: sources/stacks-more-algebra/pinned-original-source.tex, source lines 31453--31453. Same-chapter navigation pointer; the statement it names is standard for the section and the argument that consumes it stays inside the counted proof span.\par\medskip
\begin{lemma}
\label{lemma-minimal-primes-tensor-strictly-henselian}
Let $k$ be an algebraically closed field. Let $A$, $B$ be strictly
henselian local $k$-algebras with residue field equal to $k$.
Let $C$ be the strict henselization of $A \otimes_k B$ at the maximal
ideal $\mathfrak m_A \otimes_k B + A \otimes_k \mathfrak m_B$.
Then the minimal primes of $C$ correspond $1$-to-$1$ to pairs of
minimal primes of $A$ and $B$.
\end{lemma}

\begin{proof}
First note that a minimal prime $\mathfrak r$ of $C$ maps to a minimal
prime $\mathfrak p$ in $A$ and to a minimal prime $\mathfrak q$ of $B$
because the ring maps $A \to C$ and $B \to C$ are flat (by going down for
flat ring map Algebra, Lemma \ref{algebra-lemma-flat-going-down}).
Hence it suffices to show that the strict henselization of
$(A/\mathfrak p \otimes_k B/\mathfrak q)_{
\mathfrak m_A \otimes_k B + A \otimes_k \mathfrak m_B}$
has a unique minimal prime ideal. By
Algebra, Lemma \ref{algebra-lemma-quotient-strict-henselization}
the rings $A/\mathfrak p$, $B/\mathfrak q$ are strictly henselian.
Hence we may assume that $A$ and $B$ are strictly henselian
local domains and our goal is to show that $C$ has a unique minimal prime.
By Lemma \ref{lemma-geometrically-unibranch} the
integral closure $A'$ of $A$ in its fraction field
is a normal local domain with residue field $k$. Similarly for the
integral closure $B'$ of $B$ into its fraction field. By
Algebra, Lemma \ref{algebra-lemma-geometrically-normal-tensor-normal}
we see that $A' \otimes_k B'$ is a normal ring. Hence its localization
$$
R = (A' \otimes_k B')_{
\mathfrak m_{A'} \otimes_k B' + A' \otimes_k \mathfrak m_{B'}}
$$
is a normal local domain. Note that $A \otimes_k B \to A' \otimes_k B'$
is integral (hence gong up holds --
Algebra, Lemma \ref{algebra-lemma-integral-going-up})
and that $\mathfrak m_{A'} \otimes_k B' + A' \otimes_k \mathfrak m_{B'}$
is the unique maximal ideal of $A' \otimes_k B'$
lying over $\mathfrak m_A \otimes_k B + A \otimes_k \mathfrak m_B$.
Hence we see that
$$
R = (A' \otimes_k B')_{
\mathfrak m_A \otimes_k B + A \otimes_k \mathfrak m_B}
$$
by
Algebra, Lemma \ref{algebra-lemma-unique-prime-over-localize-below}.
It follows that
$$
(A \otimes_k B)_{
\mathfrak m_A \otimes_k B + A \otimes_k \mathfrak m_B}
\longrightarrow
R
$$
is integral. We conclude that $R$ is the integral closure of
$(A \otimes_k B)_{
\mathfrak m_A \otimes_k B + A \otimes_k \mathfrak m_B}$
in its fraction field, and by
Lemma \ref{lemma-geometrically-unibranch}
once again we conclude that $C$ has a unique prime ideal.
\end{proof}

\end{document}
